\documentclass[11pt,a4paper]{article}
\usepackage[utf8]{inputenc}
\usepackage[T1]{fontenc}
\usepackage[brazilian]{babel}
\usepackage[margin=2cm]{geometry}
\usepackage{lmodern}
\usepackage{amssymb}
\usepackage{xcolor}
\usepackage{array}
\usepackage{booktabs}
\usepackage{enumitem}
\usepackage{titlesec}
\usepackage{fancyhdr}
\usepackage{parskip}

\definecolor{brandgreen}{HTML}{2E5E3E}
\definecolor{lightbg}{HTML}{F5F2EC}

\titleformat{\section}{\large\bfseries\color{brandgreen}}{}{0pt}{}[{\color{brandgreen}\titlerule}]
\titleformat{\subsection}{\bfseries\color{brandgreen}}{}{0pt}{}

\pagestyle{fancy}
\fancyhf{}
\fancyhead[L]{\small Querência ERP --- Avaliação de usabilidade}
\fancyhead[R]{\small Fluxo 1 --- Diretoria}
\fancyfoot[C]{\small\thepage}
\renewcommand{\headrulewidth}{0.4pt}

\newcommand{\dizer}[1]{%
  \noindent\colorbox{lightbg}{%
    \begin{minipage}{0.95\linewidth}
    \vspace{1mm}
    \textbf{O que dizer ao participante:} ``#1''
    \vspace{1mm}
    \end{minipage}%
  }\vspace{3mm}
}
\newcommand{\criterio}[1]{\noindent\textit{Critério de sucesso objetivo:} #1}
\newcommand{\estado}[1]{\noindent\textit{Estado necessário (já preparado no kit):} #1\vspace{2mm}}

\begin{document}

\begin{center}
  {\LARGE\bfseries\color{brandgreen} Fluxo 1 --- Cadastro de associado}\\[1mm]
  {\large RF01, RF18 --- Papel: Diretoria (painel web)}
\end{center}

\vspace{2mm}
\estado{ficha de persona ainda não usada.}

\dizer{Alguém novo quer se associar à entidade. Você recebeu esta ficha preenchida à mão ---
cadastre essa pessoa no sistema como se fosse um caso real.}

\begin{enumerate}[leftmargin=6mm]
  \item Entregar a ficha de persona do kit (ainda não usada).
  \item Participante acessa \textbf{Associados → Novo associado} (\texttt{/associados/novo}).
  \item Preenche: Nome completo, CPF, Telefone, E-mail.
  \item Preenche o endereço.
  \item Seleciona a \textbf{Categoria de sócio} indicada na ficha.
  \item Clica \textbf{Salvar}.
  \item Confirma que o associado aparece na listagem (\texttt{/associados}), buscando pelo nome ou
    CPF da ficha.
\end{enumerate}

\textbf{Verificação objetiva} (rodar depois do passo 7):
\begin{verbatim}
cd backend
node scripts/verificar-tarefas.js 1 <CPF-DA-FICHA>
\end{verbatim}
\criterio{nos dois critérios de status (associado existe, \texttt{ativo}). A conta de login não
aparece vinculada --- é normal, aqui só a diretoria cadastrou.}

\vfill
\noindent\footnotesize Se a tarefa travar por mais de \textasciitilde2 min, aplicar 1 intervenção
pontual e registrar; se mesmo assim não avançar, parar e marcar como \textit{não concluída} ---
protocolo completo no \texttt{roteiro-moderador.pdf}.

\end{document}
